\documentclass{article}
\usepackage{geometry}
\geometry{margin=1.5cm, vmargin={0pt,1cm}}
\setlength{\topmargin}{-1cm}
\setlength{\headheight}{12.5pt}
\setlength{\paperheight}{29.7cm}
\setlength{\textheight}{25.3cm}

% useful packages.
\usepackage[UTF8]{ctex}
\usepackage{fancyhdr}
\usepackage{layout}
\usepackage{luacode}

% import common macros
% % % % % % % % % % % % % % % % % % % % % % % % % % % % % % % %
% Common Macros for LaTeX
% % % % % % % % % % % % % % % % % % % % % % % % % % % % % % % %

% Useful packages
\usepackage{amsfonts}
\usepackage{amsmath}
\usepackage{amssymb}
\usepackage{amsthm}
\usepackage{enumerate}
\usepackage{graphicx}
\usepackage{multicol}
\usepackage[ruled,linesnumbered]{algorithm2e}

% Math operators and common symbols
\newcommand{\dif}{\mathrm{d}}
\newcommand{\innerProd}[1]{\left\langle #1 \right\rangle}
\newcommand{\difFrac}[2]{\frac{\dif #1}{\dif #2}}
\newcommand{\pdfFrac}[2]{\frac{\partial #1}{\partial #2}}
\newcommand{\T}{\mathrm{T}}
\newcommand{\norm}[1]{\left\| #1 \right\|}
\newcommand{\abs}[1]{\left| #1 \right|}

% Vector and Matrix shortcuts
\newcommand{\vecTwo}[2]{
  \begin{bmatrix}
    #1 \\
    #2
\end{bmatrix}}
\newcommand{\vecThree}[3]{
  \begin{bmatrix}
    #1 \\
    #2 \\
    #3
\end{bmatrix}}
\newcommand{\vecTwoH}[2]{
  \begin{bmatrix}
    #1 & #2
  \end{bmatrix}
}
\newcommand{\vecThreeH}[3]{
  \begin{bmatrix}
    #1 & #2 & #3
  \end{bmatrix}
}

% Common fields
\newcommand{\R}{\mathbb{R}}
\newcommand{\C}{\mathbb{C}}
\newcommand{\Z}{\mathbb{Z}}
\newcommand{\N}{\mathbb{N}}

% Solutions and proofs
\newenvironment{solution}{
\begin{proof}[解]}{
\end{proof}}

% Utility to get environment variables (requires lualatex)
\newcommand{\getEnv}[2]{\directlua{tex.print(os.getenv("#1") or "#2")}}


\newcommand{\tridiag}{\operatorname{tridiag}}
\newcommand{\off}{\operatorname{off}}

\usepackage{listings}
\usepackage{xcolor}

% Configure listings for Python code highlighting
\lstset{
  language=Python,
  basicstyle=\ttfamily\small,
  keywordstyle=\bfseries\color{blue},
  commentstyle=\color{gray},
  stringstyle=\color{red},
  breaklines=true,
  showstringspaces=false,
  tabsize=4,
  frame=single,
  framesep=5pt,
  rulecolor=\color{gray},
  backgroundcolor=\color{white},
  captionpos=b,
  numberstyle=\tiny\color{gray},
  numbers=left,
  stepnumber=5,
}

\lstnewenvironment{plaintext}[1][]{
  \lstset{
    language=none,
    basicstyle=\ttfamily\small,
    keywordstyle=,
    commentstyle=,
    stringstyle=,
    numbers=none,
    frame=single,
    backgroundcolor=\color{white},
    #1
  }
}{}

% Name and uID
\newcommand{\name}{\getEnv{NAME}{NAME}}
\newcommand{\uid}{\getEnv{UID}{UID}}
\newcommand{\course}{数值代数}
\newcommand{\hwNum}{15}

\begin{document}

\pagestyle{fancy}
\fancyhead{}
\lhead{\name (\uid)}
\chead{\course \#\hwNum}
\rhead{\today}

\section{实对称三对角阵的过关 Jacobi 方法}

求实对称三对角阵的全部特征值和特征向量.

\begin{enumerate}[(1)]
  \item 用熟悉的计算机语言编制利用过关 Jacobi 方法求实对称三对角阵全部特征值和特征向量的通用子程序.
  \item 利用所编制的子程序求矩阵（从 $50$ 阶到 $100$ 阶）
    \[
      A = \tridiag(1,4,1)
      =
      \begin{bmatrix}
        4 & 1 \\
        1 & 4 & 1 \\
          & \ddots & \ddots & \ddots \\
          &        & 1 & 4 & 1 \\
          &        &   & 1 & 4
      \end{bmatrix}
    \]
    的全部特征值和特征向量.
\end{enumerate}

\begin{solution}

  设当前矩阵为 $B=(b_{ij})$. 过关 Jacobi 方法每一趟先取
  \[
    \eta = \frac{\norm{\off(B)}_F}{4n},
  \]
  然后按 $1\le p<q\le n$ 扫描上三角元. 只有当 $\abs{b_{pq}}\ge \eta$ 时，才对
  $(p,q)$ 平面作 Jacobi 旋转. 旋转参数取
  \[
    \tau = \frac{b_{qq}-b_{pp}}{2b_{pq}},\qquad
    t = \frac{\operatorname{sgn}(\tau)}{\abs{\tau}+\sqrt{1+\tau^2}},\qquad
    c = \frac{1}{\sqrt{1+t^2}},\quad s=tc .
  \]
  用 $J(p,q,c,s)$ 作相似变换 $B\leftarrow J^\T BJ$，同时累积
  $V\leftarrow VJ$. 当 $\norm{\off(B)}_F$ 小于给定容差时，$B$ 的对角元为全部特征值，
  $V$ 的列向量为对应特征向量.

  对本题的 Toeplitz 三对角矩阵，还可以用解析公式检查结果：
  \[
    \lambda_k = 4 + 2\cos\frac{k\pi}{n+1},\qquad
    u_j^{(k)} = \sqrt{\frac{2}{n+1}}\sin\frac{jk\pi}{n+1},
    \quad k,j=1,\ldots,n.
  \]
  程序中不使用该公式求解，只用它计算误差.

  对 $n=50,51,\ldots,100$ 的全部阶数，过关 Jacobi 程序都完成了全部特征值和特征向量计算.
  其中 $n=100$ 时：
  \[
    \lambda_{\min}=2.000967435416,\qquad
    \lambda_{\max}=5.999032564584.
  \]
  此时最大特征值误差约为 $1.33\times10^{-14}$，最大特征向量误差约为
  $1.22\times10^{-10}$，最大残差
  $\max_j\norm{Av_j-\lambda_jv_j}_2$ 约为 $4.19\times10^{-12}$，
  正交性误差 $\norm{V^\T V-I}_\infty$ 约为 $4.59\times10^{-14}$.

  运行输出中逐阶列出边端特征值、误差和残差；由于完整特征向量表过长，只打印了
  $n=50,75,100$ 的若干代表性特征向量分量. 程序本身对每个 $n$ 均返回完整的
  $\lambda_1,\ldots,\lambda_n$ 和 $v_1,\ldots,v_n$.

  代码与运行结果如下：
  \lstinputlisting[language=Python, caption={实对称三对角矩阵的过关 Jacobi 方法}]{exercise.py}
  \lstinputlisting[language={}, caption={过关 Jacobi 运行结果}]{output.txt}

\end{solution}

\end{document}
